% =====================================================================
% pertemuan-07.tex
% Pertemuan 7 - Ambiguitas, Left Recursion, dan Left Factoring
% Sub-CPMK 3.2 - Praktik terbimbing + kuis konsep
% =====================================================================
\documentclass[aspectratio=169,10pt,t]{beamer}
\usepackage{config/beamerthemeubb}

\UBBmeeting{7}

\begin{document}

\UBBtitleframe{7}{Ambiguitas, Left Recursion, dan Left Factoring}{Praktik terbimbing \textperiodcentered{} Kuis konsep}

% ---------------------------------------------------------------------
\begin{frame}{Capaian Pertemuan}
  \textbf{CPMK induk:} CPMK 03 --- Merancang tata bahasa
  (\textit{grammar}) bebas konteks untuk bahasa sederhana.

  \vspace{3pt}
  \begin{block}{Sub-CPMK yang dinilai pada pertemuan ini}
    \small
    \textbf{Sub-CPMK 3.2:} Menguji grammar subset bahasa C dan
    men-\textit{debug} ambiguitas, \textit{left recursion}, serta
    \textit{left factoring} sampai berkas grammar bebas konflik.
  \end{block}

  \vspace{3pt}
  \begin{exampleblock}{Luaran yang diperiksa}
    \footnotesize
    Berkas \file{grammar-ll1.ebnf} yang bebas konflik, log sebelum dan
    sesudah \eng{debug}, lima string uji, dan bukti kesetaraan bahasa.
  \end{exampleblock}
\end{frame}

% ---------------------------------------------------------------------
\begin{frame}{Agenda Pertemuan (150 menit)}
  \small
  \begin{tabular}{@{}p{1.4cm}p{12.3cm}@{}}
    \toprule
    \textbf{Menit} & \textbf{Kegiatan} \\
    \midrule
    25 & Ceramah: ambiguitas dan \textit{dangling else} \\
    20 & Latihan: precedence, associativity, dan dua pohon urai \\
    30 & Latihan: menghapus \textit{left recursion} \\
    30 & Latihan: \textit{left factoring} dan konflik FIRST \\
    20 & Kuis konsep dan pembahasan \\
    25 & Praktik P2-B: uji dan \textit{debug} grammar \\
    \bottomrule
  \end{tabular}

  \vspace{6pt}
  {\scriptsize\color{ubbblue!75}Porsi praktik dan latihan 105 menit;
  kuis konsep 20 menit di tengah pertemuan.}
\end{frame}

% ---------------------------------------------------------------------
\begin{frame}{Posisi Pertemuan pada Rantai Proyek}
  \UBBchain{2}

  \vspace{2pt}
  \begin{itemize}
    \item Pertemuan ini membersihkan simpul \textbf{grammar} dari
          ambiguitas, rekursi kiri, dan prefiks umum.
    \item Grammar hasil transformasi menjadi syarat parser top-down
          pada Pertemuan 9.
    \item Berkas \file{grammar-ll1.ebnf} adalah artefak Sub-CPMK 3.2.
  \end{itemize}
\end{frame}

% ---------------------------------------------------------------------
\begin{frame}{Ambiguitas Grammar}
  \begin{block}{Definisi}
    \small
    Tata bahasa \(G\) ambigu jika terdapat string \(w \in L(G)\) yang
    memiliki dua pohon urai berbeda.
  \end{block}

  \begin{itemize}
    \item Parser tidak dapat menentukan struktur yang dimaksud tanpa
          aturan tambahan.
    \item Contoh: \(E \to E+E \mid E*E \mid \code{id}\).
    \item String \(\code{id}+\code{id}*\code{id}\) memiliki dua pohon:
          satu mengutamakan \(+\) di akar, satu mengutamakan \(*\).
    \item Tanpa aturan prioritas, makna ekspresi tidak jelas.
  \end{itemize}
\end{frame}

% ---------------------------------------------------------------------
\begin{frame}{Dua Pohon Urai untuk \texttt{id + id * id}}
  \begin{columns}[T,onlytextwidth]
    \begin{column}{0.47\textwidth}
      \centering
      \begin{tikzpicture}[
          nd/.style={draw=ubbblue!40,fill=ubbblue!7,rounded corners=2pt,
            minimum width=0.75cm,minimum height=0.48cm,font=\scriptsize},
          ar/.style={-{Stealth[length=1.3mm]},draw=ubbblue!55}
        ]
        \node[nd] (r) at (0,0) {\(+\)};
        \node[nd] (a) at (-1.1,-0.8) {\code{id}};
        \node[nd] (m) at (1.1,-0.8) {\(*\)};
        \node[nd] (b) at (0.4,-1.6) {\code{id}};
        \node[nd] (c) at (1.8,-1.6) {\code{id}};
        \draw[ar] (r) -- (a); \draw[ar] (r) -- (m);
        \draw[ar] (m) -- (b); \draw[ar] (m) -- (c);
      \end{tikzpicture}
      \par {\scriptsize Pohon 1: akar \(+\)}
    \end{column}
    \begin{column}{0.47\textwidth}
      \centering
      \begin{tikzpicture}[
          nd/.style={draw=ubbgold,fill=ubbgold!16,rounded corners=2pt,
            minimum width=0.75cm,minimum height=0.48cm,font=\scriptsize},
          ar/.style={-{Stealth[length=1.3mm]},draw=ubbgold!80}
        ]
        \node[nd] (r) at (0,0) {\(*\)};
        \node[nd] (m) at (-1.1,-0.8) {\(+\)};
        \node[nd] (c) at (1.1,-0.8) {\code{id}};
        \node[nd] (a) at (-1.8,-1.6) {\code{id}};
        \node[nd] (b) at (-0.4,-1.6) {\code{id}};
        \draw[ar] (r) -- (m); \draw[ar] (r) -- (c);
        \draw[ar] (m) -- (a); \draw[ar] (m) -- (b);
      \end{tikzpicture}
      \par {\scriptsize Pohon 2: akar \(*\)}
    \end{column}
  \end{columns}

  \vspace{3pt}
  {\footnotesize Keduanya valid menurut grammar. Aturan precedence
  memilih pohon 1 karena perkalian diikat lebih kuat.}
\end{frame}

% ---------------------------------------------------------------------
\begin{frame}{Tiga Cara Menangani Ambiguitas}
  \begin{enumerate}
    \setlength{\itemsep}{4pt}
    \item \textbf{Restrukturisasi grammar}: membuat tingkat prioritas
          terpisah.
    \item \textbf{Aturan precedence dan associativity}: dipakai parser
          generator.
    \item \textbf{Aturan semantik}: memilih pohon berdasarkan konteks,
          misalnya \eng{else} terdekat.
  \end{enumerate}

  \vspace{6pt}
  \begin{block}{Pilihan paling bersih}
    \footnotesize
    Restrukturisasi grammar menghilangkan ambiguitas pada tingkat
    sintaksis. Grammar aritmetika \(E\), \(T\), \(F\) adalah contohnya:
    satu nonterminal untuk satu tingkat prioritas.
  \end{block}
\end{frame}

% ---------------------------------------------------------------------
\begin{frame}[fragile]{\textit{Dangling Else}}
  \begin{columns}[T,onlytextwidth]
    \begin{column}{0.48\textwidth}
      \begin{lstlisting}[language=C,numbers=none,
        caption={Dua kemungkinan pemasangan \texttt{else}}]
if (a)
    if (b) x;
    else y;      // else milik if(b)

if (a)
    if (b) x;
else y;          // else milik if(a)
      \end{lstlisting}
    \end{column}
    \begin{column}{0.48\textwidth}
      \begin{itemize}
        \item String: \texttt{if (a) if (b) x; else y;}.
        \item \eng{Else} dapat dipasangkan ke \texttt{if} dalam atau
              luar.
        \item Aturan hampir semua bahasa: \textbf{else terdekat}.
        \item Ditulis dengan grammar \textit{matched/unmatched} atau
              prioritas \eng{shift} pada parser.
        \item Blok eksplisit memaksa pasangan luar.
      \end{itemize}
    \end{column}
  \end{columns}
\end{frame}

% ---------------------------------------------------------------------
\begin{frame}{Precedence dan Associativity}
  \begin{columns}[T,onlytextwidth]
    \begin{column}{0.58\textwidth}
      \scriptsize
      \setlength{\tabcolsep}{4pt}
      \begin{tabular}{@{}c l l@{}}
        \toprule
        \textbf{Tingkat} & \textbf{Operator} & \textbf{Associativity} \\
        \midrule
        1 (terendah) & \texttt{||} & kiri \\
        2 & \texttt{\&\&} & kiri \\
        3 & \texttt{==}, \texttt{!=} & kiri \\
        4 & \texttt{<}, \texttt{<=}, \texttt{>}, \texttt{>=} & kiri \\
        5 & \texttt{+}, \texttt{-} & kiri \\
        6 & \texttt{*}, \texttt{/}, \texttt{\%} & kiri \\
        7 & \texttt{!}, unary \texttt{-}, unary \texttt{+} & kanan \\
        8 (tertinggi) & \texttt{()}, call, literal, identifier & --- \\
        \bottomrule
      \end{tabular}
    \end{column}
    \begin{column}{0.38\textwidth}
      \begin{itemize}
        \item Precedence: operator mana yang diikat lebih kuat.
        \item Associativity: pengelompokan pada precedence sama.
        \item Kiri: \(a-b-c=(a-b)-c\).
        \item Kanan: \(a=b=c=a=(b=c)\).
        \item Non-asosiatif: \(a<b<c\) harus memakai kurung.
      \end{itemize}
    \end{column}
  \end{columns}
\end{frame}

% ---------------------------------------------------------------------
\begin{frame}[fragile]{Menyatakan Precedence dalam Grammar}
  \begin{itemize}
    \item Satu nonterminal untuk satu tingkat prioritas.
    \item Tingkat bawah menjadi anak dari tingkat atas.
    \item Pohon otomatis mencerminkan prioritas tanpa aturan tambahan.
  \end{itemize}

  \vspace{2pt}
  \begin{columns}[T,onlytextwidth]
    \begin{column}{0.48\textwidth}
      \begin{block}{Contoh Subset-C}
        \footnotesize
        \(\text{additive} \to \text{term}\ \{(\texttt{"+"} \mid
        \texttt{"-"})\ \text{term}\}\).\\
        \texttt{term} menangani perkalian, sehingga perkalian selalu
        lebih dalam.
      \end{block}
    \end{column}
    \begin{column}{0.48\textwidth}
      \begin{lstlisting}[language=Bison,numbers=none,
        caption={Deklarasi precedence pada Bison}]
%left OR
%left AND
%left EQ NE
%left LT LE GT GE
%left PLUS MINUS
%left STAR SLASH PERCENT
%right NOT UMINUS
      \end{lstlisting}
    \end{column}
  \end{columns}
\end{frame}

% ---------------------------------------------------------------------
\begin{frame}{\textit{Left Recursion}: Masalah dan Penghapusan}
  \begin{block}{Mengapa bermasalah}
    \small
    Parser top-down memanggil fungsi untuk nonterminal. Aturan
    \(A \to A\alpha\) membuat fungsi \texttt{parseA} memanggil dirinya
    sendiri sebelum mengonsumsi token, sehingga rekursi tak berujung.
  \end{block}

  \begin{block}{Penghapusan langsung}
    \small
    \(A \to A\alpha_1 \mid A\alpha_2 \mid \beta_1 \mid \beta_2\)\\
    menjadi\\
    \(A \to \beta_1 A' \mid \beta_2 A'\),\quad
    \(A' \to \alpha_1 A' \mid \alpha_2 A' \mid \varepsilon\).
  \end{block}

  \vspace{2pt}
  {\footnotesize \(A'\) menangani pengulangan; \(A\) memilih
  alternatif dasar.}
\end{frame}

% ---------------------------------------------------------------------
\begin{frame}{Contoh: Menghapus \textit{Left Recursion}}
  \begin{exampleblock}{Ekspresi aritmetika}
    \small
    \(E \to E+T \mid T\)\\
    menjadi\\
    \(E \to T E'\),\qquad \(E' \to +T E' \mid \varepsilon\).
  \end{exampleblock}

  \begin{exampleblock}{Empat alternatif}
    \small
    \(A \to Aa \mid Ab \mid c \mid d\)\\
    menjadi\\
    \(A \to cA' \mid dA'\),\qquad
    \(A' \to aA' \mid bA' \mid \varepsilon\).
  \end{exampleblock}

  \vspace{2pt}
  {\footnotesize Pohon urai tetap sama; parser dapat memutuskan
  berdasarkan token pertama.}
\end{frame}

% ---------------------------------------------------------------------
\begin{frame}{\textit{Left Recursion} Tidak Langsung dan Kondisi Nullable}
  \begin{block}{Contoh tidak langsung}
    \footnotesize
    \(A \to Ba \mid c\), \(B \to Ab \mid d\).\\
    Substitusi \(B\) ke \(A\): \(A \to Aba \mid da \mid c\).\\
    Hapus langsung: \(A \to daA' \mid cA'\),
    \(A' \to baA' \mid \varepsilon\).
  \end{block}

  \begin{enumerate}
    \setlength{\itemsep}{1pt}
    \item Urutkan nonterminal \(A_1,\ldots,A_n\).
    \item Untuk \(A_i\), ganti aturan yang dimulai \(A_j\) (\(j<i\)).
    \item Hapus \textit{left recursion} langsung pada \(A_i\).
  \end{enumerate}

  \vspace{2pt}
  {\footnotesize\color{ubbblue!80}Jika \(\alpha\) nullable, waspadai
  loop epsilon. Periksa FIRST/FOLLOW setelah transformasi.}
\end{frame}

% ---------------------------------------------------------------------
\begin{frame}{\textit{Left Factoring}}
  \begin{block}{Transformasi}
    \small
    \(A \to \alpha\beta_1 \mid \alpha\beta_2\)\\
    menjadi\\
    \(A \to \alpha A'\),\qquad \(A' \to \beta_1 \mid \beta_2\).
    Jika salah satu alternatif berakhir setelah \(\alpha\), tambahkan
    \(\varepsilon\) pada \(A'\).
  \end{block}

  \begin{exampleblock}{If-else}
    \footnotesize
    \(S \to \texttt{if}\ E\ \texttt{then}\ S\ \texttt{else}\ S
     \mid \texttt{if}\ E\ \texttt{then}\ S\)\\
    menjadi
    \(S \to \texttt{if}\ E\ \texttt{then}\ S\ S'\),
    \(S' \to \texttt{else}\ S \mid \varepsilon\).
  \end{exampleblock}

  \vspace{2pt}
  {\footnotesize Parser memutuskan setelah selesai memproses \(S\);
  jika token berikutnya \texttt{else}, pilih alternatif pertama.}
\end{frame}

% ---------------------------------------------------------------------
\begin{frame}[fragile]{Contoh: Identifier dan Pemanggilan}
  \begin{block}{Sebelum}
    \small
    \(\text{primary} \to \texttt{ID} \mid \texttt{ID}\
    \texttt{"("}\ \text{args}\ \texttt{")"}\).
  \end{block}

  \begin{block}{Sesudah}
    \small
    \(\text{primary} \to \texttt{ID}\ \text{primary}'\),\quad
    \(\text{primary}' \to \texttt{"("}\ \text{args}\
    \texttt{")"} \mid \varepsilon\).
  \end{block}

  \vspace{3pt}
  \begin{itemize}
    \item Lexer membaca \texttt{ID} sekali.
    \item Parser memutuskan berdasarkan token berikutnya.
    \item Faktor umum pada Subset-C juga muncul pada \texttt{stmt}
          (assignment versus call).
  \end{itemize}
\end{frame}

% ---------------------------------------------------------------------
\begin{frame}{Transformasi Grammar Subset-C}
  \scriptsize
  \setlength{\tabcolsep}{4pt}
  \begin{tabular}{@{}p{2.4cm}p{6.2cm}p{4.8cm}@{}}
    \toprule
    \textbf{Aturan} & \textbf{Perubahan} & \textbf{Alasan} \\
    \midrule
    \texttt{primary} & \texttt{ID [ "(" [args] ")" ]} &
      menghilangkan prefiks \texttt{ID} \\
    \texttt{stmt} & \texttt{ID assign\_or\_call ";"} &
      memisahkan assignment dan call \\
    \texttt{expr} & tanpa assignment &
      menghindari konflik FIRST \\
    \texttt{ifstmt} & \texttt{else} opsional &
      \textit{dangling else} diselesaikan parser \\
    \bottomrule
  \end{tabular}

  \vspace{5pt}
  {\scriptsize\color{ubbblue!75}Setelah factoring, grammar lebih cocok
  untuk parser LL(1). Grammar asli tetap disimpan untuk dokumentasi.}
\end{frame}

% ---------------------------------------------------------------------
\begin{frame}[fragile]{Berkas \file{grammar-ll1.ebnf}}
  \begin{columns}[T,onlytextwidth]
    \begin{column}{0.58\textwidth}
      \lstinputlisting[language=EBNF,firstline=4,lastline=17,
        firstnumber=4,
        caption={Bagian awal \file{grammar-ll1.ebnf}},
        label={lst:p7-ll1}]{../buku_ajar/code/subsetc/spec/grammar-ll1.ebnf}
    \end{column}
    \begin{column}{0.38\textwidth}
      \begin{itemize}
        \item \texttt{assign\_or\_call} menangani \texttt{= expr} atau
              \texttt{( [args] )}.
        \item \texttt{stmt} memakai satu prefiks \texttt{ID}.
        \item \texttt{expr} tidak lagi memuat assignment.
        \item Tingkat precedence dipertahankan.
      \end{itemize}
    \end{column}
  \end{columns}
\end{frame}

% ---------------------------------------------------------------------
\begin{frame}{Alur Uji dan \textit{Debug} Grammar}
  \scriptsize
  \setlength{\tabcolsep}{4pt}
  \begin{tabular}{@{}p{4.0cm}p{4.2cm}p{5.2cm}@{}}
    \toprule
    \textbf{Gejala} & \textbf{Penyebab} & \textbf{Tindakan} \\
    \midrule
    Dua pohon urai untuk satu string & Ambiguitas &
      Restrukturisasi atau aturan resolusi \\
    Parser tidak pernah mengonsumsi token & \textit{Left recursion} &
      Hapus langsung atau tidak langsung \\
    Dua alternatif berebut token pertama & Prefiks umum &
      \textit{Left factoring} \\
    Sel ganda pada tabel LL(1) & Konflik FIRST/FOLLOW &
      Factoring, hapus rekursi, atau tambah lookahead \\
    \texttt{else} masuk dua \texttt{if} & \textit{Dangling else} &
      Aturan else terdekat atau grammar matched/unmatched \\
    \bottomrule
  \end{tabular}

  \vspace{4pt}
  {\scriptsize Uji lima string + satu kasus dangling else; terapkan
  satu transformasi pada satu waktu; uji ulang string yang sama.}
\end{frame}

% ---------------------------------------------------------------------
\begin{frame}{Kuis Konsep}
  \small
  \begin{enumerate}
    \setlength{\itemsep}{3pt}
    \item Apa definisi tata bahasa ambigu?
    \item Mengapa \textit{left recursion} menghambat parser top-down?
    \item Ubah \(A \to Aa \mid b\) menjadi bentuk tanpa
          \textit{left recursion}.
    \item Kapan \textit{left factoring} diperlukan?
    \item Apa aturan resolusi untuk \textit{dangling else}?
  \end{enumerate}

  \vspace{4pt}
  {\footnotesize\color{ubbblue!75}Kerjakan mandiri selama 8 menit,
  lalu bahas bersama. Kunci tersedia pada versi handout.}
\end{frame}

\mode<handout>{
\begin{frame}{Kunci Kuis Konsep}
  \footnotesize
  \begin{enumerate}
    \setlength{\itemsep}{2pt}
    \item Tata bahasa ambigu jika ada string yang memiliki dua pohon
          urai berbeda.
    \item Fungsi parser memanggil dirinya sendiri sebelum mengonsumsi
          token, sehingga rekursi tak berujung.
    \item \(A \to bA'\), \(A' \to aA' \mid \varepsilon\).
    \item Jika dua alternatif memiliki prefiks umum yang tidak dapat
          dibedakan dengan satu token lookahead.
    \item \textit{Else} dipasangkan ke \texttt{if} terdekat yang belum
          memiliki \texttt{else}.
  \end{enumerate}
\end{frame}
}

% ---------------------------------------------------------------------
\begin{frame}{Kegiatan Mahasiswa di Kelas: P2-B}
  \begin{block}{Uji dan \textit{debug} grammar}
    \footnotesize
    \begin{enumerate}
      \setlength{\itemsep}{1pt}
      \item Mulai dari \file{spec/grammar.ebnf}.
      \item Uji lima string, satu kasus dangling else, dan satu
            ekspresi campuran; catat konflik.
      \item Debug satu masalah pada satu waktu: ambiguitas, rekursi
            kiri, lalu factoring.
      \item Tulis hasil ke \file{spec/grammar-ll1.ebnf}.
      \item Uji ulang string yang sama; buktikan kesetaraan bahasa.
      \item Kumpulkan log sebelum dan sesudah debug.
    \end{enumerate}
  \end{block}
\end{frame}

% ---------------------------------------------------------------------
\begin{frame}{Tugas dan Artefak}
  \begin{block}{Artefak Sub-CPMK 3.2}
    Berkas grammar yang sudah dibebaskan dari ambiguitas,
    \textit{left recursion}, dan kebutuhan \textit{left factoring},
    disertai log sebelum/sesudah dan bukti uji.
  \end{block}

  \small
  \begin{itemize}
    \item Tidak ada \textit{left recursion} langsung maupun tidak
          langsung.
    \item Prefiks umum sudah difaktorkan; konflik FIRST hilang.
    \item \textit{Dangling else} memiliki aturan resolusi eksplisit.
    \item Lima string uji tetap valid; kasus dangling else benar.
    \item \file{grammar-ll1.ebnf} dan log sebelum/sesudah dikumpulkan.
  \end{itemize}

  \vspace{3pt}
  {\footnotesize\color{ubbblue!75}P2-B menjadi bagian Praktikum 2
  (12\%) yang dilaporkan pada Pertemuan 9.}
\end{frame}

% ---------------------------------------------------------------------
\begin{frame}{Rangkuman dan Jembatan}
  \begin{itemize}
    \item Ambiguitas berarti satu string memiliki dua pohon urai;
          diselesaikan dengan restrukturisasi, precedence, atau aturan
          semantik.
    \item Precedence dan associativity ditentukan oleh tingkat
          nonterminal dan urutan aturan.
    \item \textit{Left recursion} dihapus dengan nonterminal bantu
          \(A'\) yang menangani pengulangan.
    \item \textit{Left factoring} mengeluarkan prefiks umum agar parser
          dapat memutuskan dengan satu token lookahead.
    \item Grammar bebas konflik dibuktikan dengan lima string uji,
          kasus dangling else, dan log sebelum/sesudah.
  \end{itemize}

  \vspace{4pt}
  \begin{exampleblock}{Pertemuan berikutnya}
    Pertemuan 9 memakai grammar yang bersih untuk menghitung FIRST dan
    FOLLOW, memutuskan LL(1), dan membangun recursive descent parser.
  \end{exampleblock}
\end{frame}

\end{document}
